\chapter{Неравновесное рассогласование и рождение проекторной пары}

Каноническое положительное действие стабилизирует замкнутый обменный мост.
Это исключило простую статическую неустойчивость, но не исключило
неравновесное рождение дефектов. Настоящая проверка выясняет раннюю интуицию
проекта: первичный толчок не создаёт вибрации из ничего, а нарушает точное
гашение сопряжённых колебаний быстрее, чем вакуум успевает восстановить
согласование.

\section{Ранний постулат и его строгий перевод}

В исходном популярном слое проекта волна, возвращающаяся в противофазе,
гасит сама себя, тогда как согласованный резонанс выживает; см.
\path{архив-2025-2026/2025-12-истоки/habr/print2.md}. Там же вакуум описан
как постоянно дрожащая геометрически фрустрированная ткань. Эти образы не
являются доказательством, но теперь допускают точный перевод:
\begin{equation}
 \text{замкнутый мост}
 \longleftrightarrow
 \text{согласованное взаимное гашение},
\end{equation}
а быстрый космологический переход можно проверить как быстрое изменение параметра
порядка с конечным временем релаксации.

Механизм Киббла связывает образование дефектов с независимо выбранными
вакуумными доменами; Зурек уточняет, что характерный размер домена задаётся
конкуренцией скорости переключения и времени релаксации. См. первичные работы
Киббла, \path{J.Phys.A 9 (1976) 1387}, Зурека,
\path{Nature 317 (1985) 505}, и лабораторную реализацию на жидких
кристаллах, \path{Science 251 (1991) 1336}.

\section{Почему одной окружностной фазы недостаточно}

Разрешим формально комплексную амплитуду моста
\begin{equation}
 c=\lambda e^{i\theta}.
\end{equation}
Точные произведения $T_c^*T_c$ и $T_cT_c^*$ зависят только от
$|c|=\lambda$. Поэтому каноническая статическая энергия не видит
$\theta$. Однако фиксированное обменное Real-условие требует
\begin{equation}
 c=\overline c,
\end{equation}
то есть не допускает произвольную фазу как новое независимое локальное
поле. Фаза может иметь смысл только как голономия уже существующей
связности или как изменение самого обменного отождествления.

Даже если допустить физическую фазовую голономию, её вакуумное пространство
есть $S^1$. В трёх пространственных измерениях
\begin{equation}
 \pi_1(S^1)=\mathbb Z,
 \qquad
 \pi_2(S^1)=0.
\end{equation}
Она поддерживает вихревые нити или потоки, но не топологически устойчивые
точечные частицы.

\begin{nogocriterion}[Запрет фазовой подмены частицы]
Быстрое переключение окружностной фазы нельзя считать механизмом рождения точечных частиц:
его вторая гомотопическая группа тривиальна. Нельзя также возвращать
удалённую Real-условием фазу как свободный скаляр без вывода из физической
голономии.
\end{nogocriterion}

\section{Правильная переменная рассогласования}

В Томе V уже построена проекторная ось
\begin{equation}
 P=nn^T,
 \qquad n\sim-n,
\end{equation}
со значениями в
\begin{equation}
 \mathcal M_{\rm ord}=\mathbb{RP}^2.
\end{equation}
Для неё
\begin{equation}
 \pi_1(\mathbb{RP}^2)=\mathbb Z_2,
 \qquad
 \pi_2(\mathbb{RP}^2)=\mathbb Z.
\end{equation}
Следовательно, быстрое переключение проекторной ориентации в трёх измерениях способно
оставить точечные ежи.

Доказанный мост спинорного накрытия даёт двум подъёмам минимального ежа степени
$+1$ и $-1$, хопфовы линии $L$ и $L^*$ и после коэффициентного умножения
локальные ориентированные классы
\begin{equation}
 +15,qquad-15.
\end{equation}
Глобальная сумма пары равна нулю. Поэтому быстрое проекторное переключение имеет
ровно нужный кинематический тип: он не обязан менять общий ориентированный
заряд, но может проявить локальную Real-пару дефектов.

\begin{theorem}[Кинематический мост от быстрого переключения к классу пятнадцать]
Среди проверенных переменных только проекторная ориентация
$\mathbb{RP}^2$ поддерживает точечные дефекты в трёх измерениях и уже имеет
доказанное проектное отображение минимальных зарядов $\pm1$ в классы
$\pm15$. Простая окружностная фаза этого свойства не имеет.
\end{theorem}

\section{Что даёт и чего не даёт механизм Киббла--Зурека}

Если около перехода
\begin{equation}
 \xi=\xi_0|\epsilon|^{-\nu},
 \qquad
 \tau=\tau_0|\epsilon|^{-z\nu},
 \qquad
 \epsilon=t/\tau_Q,
\end{equation}
то условие потери адиабатичности даёт
\begin{equation}
 \widehat\xi
 =\xi_0\left(\frac{\tau_Q}{\tau_0}\right)^{
 \nu/(1+z\nu)}.
 \label{eq:v6-kz-freezeout-length}
\end{equation}
Для точечных дефектов условная плотность масштабируется как
\begin{equation}
 n_{\rm point}\sim\widehat\xi^{-3}.
 \label{eq:v6-kz-point-density}
\end{equation}

Эти формулы доказывают структуру зависимости, но не дают проектного числа.
Текущий родитель ещё не вывел $\xi_0$, $\tau_0$, $\nu$, $z$ и закон
космологического переключения. Не доказано также, что вся вакуумная орбита
$\mathbb{RP}^2$ входит в один временной функционал с обменной амплитудой
$\lambda$.

\section{Новый статус вопроса}

Статическая неустойчивость замкнутого вакуума больше не является
единственным возможным маршрутом. Материя могла возникнуть как
неравновесный остаток быстрого перехода: вакуум локально стремится обратно
к замыканию, но топологический заряд мешает одновременно согласовать все
области. Это точная современная форма ранней идеи о несовпавших вибрациях.

Однако пока это условный механизм, а не доказанная космология. Следующий
проверка должна вывести один динамический функционал для
$(\lambda,P)$, определить его время релаксации и проверить, сохраняется ли
еж после завершения переключения без ручной стабилизации.

\begin{protocol}
\begin{enumerate}
 \item ранняя идея взаимного гашения: \textbf{найдена в архиве};
 \item статическая энергия зависит от фазы моста: \textbf{нет};
 \item произвольная фаза совместима с фиксированным Real-обменом:
 \textbf{нет};
 \item $S^1$ поддерживает точечные дефекты в трёх измерениях:
 \textbf{нет};
 \item $\mathbb{RP}^2$ поддерживает точечные ежи: \textbf{да};
 \item отображение минимального ежа в $\pm15$: \textbf{уже доказано};
 \item глобальная нейтральность пары: \textbf{сохранена};
 \item механизм быстрого рассогласования: \textbf{кинематически открыт};
 \item критические показатели и микроскопические масштабы:
 \textbf{не выведены};
 \item рождение материи: \textbf{не доказано};
 \item следующая проверка:
 \textbf{родительская динамика проекторного переключения}.
\end{enumerate}
\end{protocol}

Машинный сертификат сохранён в
\path{s2t_v6_closed_bridge_destabilization_gate_results.json}.